\section{Results}\label{sec:results}

Five subsections follow: the information present in the inputs before any
network is trained, the accuracy of the forecasts under the two protocols, the
ratio between them, the comparison with benchmarks, and the comparison between
the two heads.

\subsection{Information about the next day in the inputs}
Before any network is trained, the inputs themselves already show how the two
protocols differ. Table~\ref{tab:diagnostic} relates the high frequency group,
the sum of the $K$ VMD modes on the last day of each input window, to the price
change that the forecast has to predict. Under full series decomposition the
two are strongly and negatively correlated on every series, with $r$ between
$-0.58$ and $-0.53$ ($p<10^{-35}$), and a linear regression on three input
features, fitted on half of the test dates and scored on the other half,
removes between $51\%$ and $60\%$ of the squared error of persistence. Daily
index changes should not allow anything of the kind. Fig.~\ref{fig:mechanism}
shows why it happens: when the full series is decomposed, the high frequency
value at day $t-1$ is shaped by the price at day $t$, so a rise on the next day
leaves day $t-1$ looking like a local trough and pulls its high frequency value
down. Under walk forward decomposition the same correlation lies between
$0.00$ and $0.10$, and the regression does worse than persistence on all four
series, by $0.9\%$ to $3.0\%$. The one correlation that reaches nominal
significance, $r=0.10$ for the SSEC highest price ($p=0.03$), carries no value
out of sample. The walk forward high frequency group is also $2.5$ to $3.3$
times more variable on the last day of the history than the full series group
on the same dates, which is the boundary behaviour that sifting is known to
need care with~\cite{emdtutorial2023}. Those large end values are noise rather
than signal, and a forecaster that works from history alone has to live with
them.

%%TABLE:tab_diagnostic%%

\begin{figure}[ht]
\centering
\figorbox{\textwidth}{figs/fig_mechanism.png}
\caption{Next day change of the S\&P~500 highest price plotted against the
high frequency group on the last day of the input window, for the $496$ test
dates both protocols share. Left, full series decomposition; right, walk
forward decomposition. The figure shows the shape of the relation that
Table~\ref{tab:diagnostic} summarises with one coefficient per series.}
\label{fig:mechanism}
\end{figure}

\subsection{Forecast accuracy under the two protocols}
Table~\ref{tab:accuracy} reports the accuracy of both heads under both
protocols, with persistence alongside, on the $496$ (S\&P~500) and $478$
(SSEC) test dates that the protocols share. Under full series decomposition
the forecasts are very accurate by any of the usual measures: RMSE between
$5.0$ and $10.7$ index points, MAPE between $0.066\%$ and $0.134\%$, and $R^2$
between $0.9992$ and $0.9999$. These values are typical of what decomposition
ensemble studies report. Under walk forward decomposition the same networks,
trained with the same code, seeds and budgets, reach RMSE between $42.5$ and
$68.5$ points, MAPE between $0.72\%$ and $0.83\%$, and $R^2$ between $0.981$
and $0.991$. Persistence sits between the two, with RMSE from $36.2$ to $57.9$
points and MAPE from $0.52\%$ to $0.67\%$. Fig.~\ref{fig:trace} shows the
first hundred test days of two series, where the difference is visible without
any statistic: the full series forecast lies on top of the price and turns
when the price turns, which is only possible because its inputs were computed
from the price it is forecasting, while the walk forward forecast follows the
price with a delay, as persistence does, and with additional error of its own.

%%TABLE:tab_accuracy%%

\begin{figure}[ht]
\centering
\figorbox{\textwidth}{figs/fig_trace.png}
\caption{Price, persistence and the two TCN forecasts over the first $100$
shared test days of the S\&P~500 highest price (top) and the SSEC highest price
(bottom). The figure adds the timing of the forecasts, which the error
measures in Table~\ref{tab:accuracy} do not show.}
\label{fig:trace}
\end{figure}

\subsection{How much of the accuracy depends on the protocol}
Table~\ref{tab:leakage} and Fig.~\ref{fig:rho} give the leakage inflation
ratio of Eq.~\eqref{eq:rho} for each series and head. Restricting the
decomposition to the past multiplies the RMSE by a factor between $4.9$ and
$10.4$, with a geometric mean of $7.2$ over the eight cases. The MDM statistic
of Eq.~\eqref{eq:mdm} for the walk forward forecast against the full series
forecast lies between $+4.5$ and $+12.4$, so every one of the eight
differences is significant at any conventional level. The ratio is larger for
the TCN head than for the TRM head on every series ($6.4$ to $10.4$ against
$4.9$ to $8.3$), because the TCN head achieved the lower full series error
while the two heads reach almost the same walk forward error.

%%TABLE:tab_leakage%%

\begin{figure}[ht]
\centering
\figorbox{0.8\textwidth}{figs/fig_rho.png}
\caption{Leakage inflation ratio $\rho$ by series and head. A star marks a
significant MDM test of the walk forward forecast against the full series
forecast ($p<0.05$); the line at $\rho=1$ is where decomposing the full series
would give no advantage. The figure makes the consistency of the effect across
series and heads visible at a glance, which the rows of Table~\ref{tab:leakage}
do not.}
\label{fig:rho}
\end{figure}

\subsection{Comparison with benchmarks}
Table~\ref{tab:baselines} gives the RMSE of the five benchmarks. Persistence
and the drift forecast are indistinguishable, within $0.3\%$ of each other on
every series, and both are far ahead of any smoothing benchmark: the five day
moving average is $1.5$ to $1.6$ times worse, exponential smoothing $1.8$ to
$1.9$ times, and the twenty day moving average $2.7$ to $2.9$ times. On daily
index levels, the last observed price is the benchmark to beat.
Table~\ref{tab:dm} tests each model against it. Under full series
decomposition both heads beat persistence on all four series, with MDM
statistics between $-4.5$ and $-10.1$ and $p<0.001$ in every case. Under walk
forward decomposition both heads lose to persistence on all four series, with
MDM statistics between $+3.5$ and $+7.4$, again $p<0.001$ in every case. In
terms of RMSE, the causal forecasts are between $13\%$ and $30\%$ worse than
persistence. Of the eight model and series combinations, all eight beat
persistence when the decomposition sees the future and none does when it does
not.

%%TABLE:tab_baselines%%

%%TABLE:tab_dm%%

\subsection{TCN and TRM heads}
Table~\ref{tab:heads} compares the two heads under each protocol, using the
MDM statistic of the TCN forecast against the TRM forecast. Under full series
decomposition the TCN head is significantly better on all four series (MDM
between $-2.9$ and $-6.1$, $p<0.01$), which is the ranking the source
comparison reported. Under walk forward decomposition the picture changes: on
three series the difference is not significant ($p=0.56$, $0.52$ and $0.97$),
and on the fourth, the S\&P~500 lowest price, the TRM head is significantly
better (MDM $+4.2$, $p<0.001$). The two heads reach walk forward RMSE values
that differ by less than $0.7\%$ on three of the four series.

%%TABLE:tab_heads%%

\section{Discussion}\label{sec:discussion}

The discussion has five parts: what the full series numbers measure, why the
causal forecasts trail persistence, what the head comparison shows, the
practices we recommend, and the limits of the evidence.

\subsection{What the full series numbers measure}
The full series results in Table~\ref{tab:accuracy} are of the same order as
those reported for this family of models: percentage errors around one tenth
of one percent and coefficients of determination above $0.999$. They were
produced by a pipeline whose implementation we audited line by line, with no
test data in any scaler, training set or stopping rule. They are nonetheless
not forecasts. Table~\ref{tab:diagnostic} shows that the inputs alone, with no
network, explain half of the next day's price change under this protocol, and
Table~\ref{tab:leakage} shows that the whole advantage over persistence
disappears when the decomposition is restricted to the past. The number that a
full series protocol measures is how efficiently a network can read the future
out of components that already contain it.

\subsection{Why the causal forecasts trail persistence}
The walk forward forecasts were not merely no better than persistence; they
were significantly worse on every series, by $13\%$ to $30\%$ in RMSE. Two
measurements locate the cause. First, the components at the newest end of a
freshly decomposed history are $2.5$ to $3.3$ times more variable than the same
components inside a full decomposition, which is the boundary behaviour of
sifting~\cite{emdtutorial2023}. Second, a forecaster that simply copied the
last value of each component window and summed the copies would have matched
persistence to within $0.7\%$ under either protocol, because the components
reconstruct the price. The trained networks therefore lost between $12\%$ and
$29\%$ relative to copying their own inputs. They were trained on windows from
the interior of the training decomposition, where component dynamics are
smooth, and applied to windows that end at a boundary, where they are not. This
is a mismatch between training and test inputs that the walk forward protocol
exposes rather than creates, and it applies equally to the source model and to
the many published variants that decompose the training portion once. Closing
it would require training on windows built the same way as the test windows,
at the cost of one extra decomposition per training day, roughly two thousand
per series, or an explicit boundary treatment inside the decomposition.

\subsection{The head comparison}
Under full series decomposition the TCN head was significantly more accurate
than the TRM head on all four series, which is the conclusion the source
comparison reached. Under walk forward decomposition the difference vanished
on three series and reversed on the fourth. The ranking of the two heads is
therefore a property of the evaluation protocol, not of the heads: when the
inputs contain the answer, the architecture that reads it out more efficiently
wins, and when they do not, the two are equivalent. This is consistent with the
observation that simple models hold their own against elaborate architectures
when forecasts are evaluated carefully~\cite{zeng2023,hewamalage2022}.

\subsection{Recommendations}
Three practices would have made the source result, and results like it,
interpretable.
\begin{enumerate}
\item Decompose only data that would have been available at the forecast
origin, and say so. Table~\ref{tab:diagnostic} gives a check that takes a few
seconds and needs no model.
\item Report persistence and the MDM statistic against it, since on daily
index levels it is the benchmark that matters.
\item Report the ratio $\rho$ of causal to noncausal error, so that readers
can see how much of a headline number depends on the protocol.
\end{enumerate}

\subsection{Limitations}
Five choices bound these conclusions.
\begin{enumerate}
\item The noise amplitude and the swarm objective were kept from the source
model, so the decomposition studied is ICEEMDAN at a small, fixed noise level
with $K$ and $\alpha$ at the bounds of their search ranges. Other settings
could change the size of the effect, although not the fact that a
decomposition of the full series has seen the prices it is asked to forecast.
\item The forecast horizon is one day, and the evidence covers two indices
over one ten year window with one seed per network.
\item Early stopping monitors the training loss, which keeps every network
away from test data but may leave individual components over or under
trained.
\item Input windows are expressed relative to their last value rather than
min max scaled, a change the S\&P~500 required but one that departs from the
source model, so our absolute errors are not directly comparable with the
errors reported there.
\item Walk forward training windows come from the interior of the training
decomposition while walk forward test windows end at a boundary, as discussed
above. The causal results therefore describe the protocol as it is commonly
run, and whether a boundary matched training scheme would narrow the gap to
persistence is untested here.
\end{enumerate}
